\section{Translation}
\label{sec:translation}

The calculus presented in \cref{sec:expr-proc,sec:types,sec:semantics}
records ownership and ordering in a linear type system and, dynamically, in
binder groups.  This section defines an untyped run-time calculus that is
closer to an actual implementation of splitting and borrowing.  The
run-time calculus provides the tools, ordinary channels
and synchronization cells, so that translated well-typed source processes
preserve the ordering guarantees established by the type system and the direct
semantics of \cref{sec:semantics}.

\begin{figure}
  \setlength{\abovedisplayskip}{0pt}%
  \setlength{\belowdisplayskip}{0pt}%
  \begin{align*}
    \mathsf{Flag} &\grmeq \FlagInit \grmor \FlagDropped \\
    \mathsf{Channel} &\grmeq
      \{\StateOpen,\StateClosed\}
      \times \mathsf{Flag}^{*} \times \mathsf{Flag}^{*} \\
    \mathsf{Config} &\grmeq (C,S,E)
  \end{align*}
  \caption{Untyped run-time syntax}
  \label{fig:unr-syntax}
\end{figure}

\Cref{fig:unr-syntax} defines an untyped run-time configuration $(C,S,E)$.
Here $C$ is a finite set of primitive channel names, $S:C\to\mathsf{Channel}$
maps each channel to its current state, and $E$ is a finite multiset of untyped
expressions executing concurrently.  A channel state records whether the
channel is open and a sequence of synchronization flags for each of its two
endpoints.  A flag in state $\FlagInit$ records that the corresponding $\CDrop$
is still pending; after the drop, its state is $\FlagDropped$ and an $\CAcq$
may proceed.  The term forms for constants and ordinary expressions are as
before; the target remote-split instruction remains unary and always
allocates its fresh flag at the front of the endpoint's list.
The multiset $E$ represents parallel execution.

We call the channels of the target calculus \emph{primitive channels} to
distinguish them from (translated) source channels.  A translated source channel
is represented by a value $\targettermcolor{c}=\USHandle[\U]{\TEVary}[\V]$, a
triple consisting of an optional pre-synchronization reference $\U$, a primitive
channel endpoint $\TEVary$, and an optional post-synchronization reference $\V$.
The default optional component is $\USUnit$.  A nontrivial reference identifies
a flag in the list belonging to $\TEVary$.  Its pre-flag is consumed by $\CAcq$
before the endpoint is used; its post-flag is signaled by $\CDrop$ when
ownership of the endpoint prefix is released.  Communication, choice, and close
operate on the primitive channel; the flag references encode the ordering
constraints on remote borrows that source binder groups enforce directly. Local
borrows do not use this machinery.

Permuting concurrent expressions or channel entries does not change a
configuration: $E$ is a multiset, $C$ is a set, and $S$ is keyed by channel
name.  There are no channel or synchronization binders whose scopes need to
be rearranged.  Up to consistent renaming of fresh names, structural
congruence boils down to equality of configurations, and reduction needs no
separate congruence rule.

For $S(\TEc)=(b,\TFlags[0],\TFlags[1])$, define $\USLive=b$
and $\USFlags=\TFlags[\TEs]$, where
$\TEs\in\{\targetendpointcolor{0},\targetendpointcolor{1}\}$ selects one of the
two endpoints $\TEcs$ and $\TEs*=1-\TEs$ selects the other.  The notation
$\USSideUpdate{\TFlags}$ updates only that endpoint's flag sequence.  For
the rules, we give each occurrence in a flag sequence a distinct name.  In a
well-formed handle, its post-component $\V$ is either $\USUnit$ or a reference
to a flag on the same endpoint.  Remote split allocates a fresh flag at the
front of that endpoint's list.  The pre- and post-references in the resulting
handles encode the dependency; physical flag-list order is an addressing
convention, rather than source session order.  In the positional Agda
representation, all existing references on the endpoint are shifted consistently
when its front cell is inserted.  The instruction itself has no position
argument.

\begin{figure}
  \small
  \begin{mathpar}
    \RuleSoupRedExp \and
    \RuleSoupRedFork \and
    \RuleSoupRedDiscard \and
    \RuleSoupRedNew \and
    \RuleSoupRedLSplit \and
    \RuleSoupRedRSplit \and
    \RuleSoupRedDrop \and
    \RuleSoupRedAcquire \and
    \RuleSoupRedClose \and
    \RuleSoupRedCom \and
    \RuleSoupRedChoice
  \end{mathpar}
  \caption{Reduction of translated configurations
    ($(C,S,E) \TPReduce (C',S',E')$)}
  \label{fig:unr-process-reduction}
\end{figure}

\Cref{fig:unr-process-reduction} gives the reduction relation on configurations.
The pattern $E\uplus\USMulti{\E}$ selects one occurrence of $\E$ from the
multiset; the communication, choice, and close rules select two distinct
occurrences.  We write $\ECapp\E$ for expression $\E$ in evaluation context
$\EC$.  {\color{red}The meta variable $\V$ ranges over syntactic values, so no separate
value premise is needed.}  The remote-split rule allocates its new cell at the
front of the selected endpoint's flag list.  In the drop and acquire rules, the
name $\TEVarz$ in the selected handle identifies a unique cell in $\USFlags$;
the displayed sequence decomposition therefore fixes the cell to update or
remove.

The expression step and discard leave the channel state unchanged; fork adds
one expression to the soup.  New adds an open channel with an already-acquired
flag on each endpoint.  A local split shares its endpoint between the two
resulting handles, whereas a remote split inserts a fresh $\FlagInit$ flag at
the front of the endpoint's list.  Drop changes that flag to
$\FlagDropped$; acquire removes it once it is ready.  Communication and choice
act on opposite endpoints of an open channel.  Close requires both flag
sequences to be empty and marks the channel closed.  These steps operate
directly on the selected expressions and channel state, so no parallel or
binder closure rules are needed.

\begin{figure}
  \small
  \begin{mathpar}
    \RuleSoupTransExp \and
    \RuleSoupTransPar \and
    \RuleSoupTransNew \\
    \begin{array}{c}
      \RuleSoupTransBindNil \\[.8em]
      \RuleSoupTransBindOne \\[.8em]
      \RuleSoupTransSeqNil  \\[.8em]
      \RuleSoupTransSeqOne
    \end{array}\and
    \begin{array}{c}
      \RuleSoupTransBindSep \\[1.5em]
      \RuleSoupTransSeqMore \\[1em]
    \end{array}
  \end{mathpar}
  \smallskip
  \[
    \BindFlag{\Bind} = \FlagDropped \text{ if } \Bind = \BEmp;
    \quad \BindFlag{\Bind} = \FlagInit \text{ otherwise.}
  \]
  \caption{Translation of processes ($\USTransProc{\Proc}=(C,S,E)$), binder
    groups ($\USTransBind{\BindGroup}{\TEcs}{\U}{\V}=(\Sub,\TFlags)$), and
    sequential segments ($\USTransSeq{\Bind}{\TEcs}{\U}{\V}=\Sub$).}
  \label{fig:translation}
\end{figure}

\Cref{fig:translation} defines a configuration-valued translation
$\USTransProc{\Proc}$.  Its environment $\Sub$ maps free source variables to
target expressions, with plain channel triples for endpoint variables. An
expression process contributes an to the multiset. Other than applying the
substitution $\Sub$ to it, no translation of expressions occurs. Parallel
composition combines disjoint channel sets and maps and takes the multiset union
of expressions.  A restriction chooses a fresh primitive channel $\TEc$,
translates its two binder groups for endpoints $\TEc[0]$ and $\TEc[1]$, and
translates the body under the substitutions they produce.  It records the
resulting flag sequences in $S(\TEc)$ and adds no expression of its own.

The binder-group function $\mathcal B$ returns a substitution whose domain is
exactly the bound variables, together with an ordered flag sequence.  Its helper
$\mathcal L$ distributes a plain channel triple across one sequential segment:
the first variable receives the pre-reference, the last the post-reference, and
intermediate variables receive neither.  A sole variable receives both.  Each
separator introduces a fresh flag name $\TEVarz$; the left segment uses
$\TEVarz$ as its post-reference and the remaining group uses it as its
pre-reference.  The flags initially occur in separator order.
$\BindFlag{\Bind}$ initializes a flag to $\FlagDropped$ when its left segment is
empty, and to $\FlagInit$ otherwise.  Substitution unions have disjoint domains:
restriction extends the incoming substitution with the bindings for its two
endpoints before translating the body.

The initial flag lists follow source separator order, whereas remote split
allocates each new cell at the front.  Synchronization depends on the
pre- and post-references in the handles, so physical storage order can vary
while preserving the same dependencies.  In the positional representation
used by Agda, front insertion shifts every existing reference on that
endpoint by one throughout the expression multiset, including inside
closures and evaluation contexts.  Acquire removes its referenced cell and
renumbers the surviving references.  With the named references used in
\cref{fig:unr-process-reduction}, both updates preserve the names of surviving
cells.  The image relation in \cref{sec:simulation} accounts for these
storage choices by consistent flag permutations.

%%% Local Variables:
%%% mode: latex
%%% TeX-master: "../popl2027.tex"
%%% End:
